\chapter{Result Format}

FEMaster writes two result formats by default. The native \file{.res} file is a text format that preserves all fields written by the loadcases. The \file{.frd} file is an ASCII CalculiX/CGX-style result file for supported nodal result visualization.

Both files use the same job base name. For example, an input deck \file{model.inp} produces \file{model.res} and \file{model.frd} unless \key{--output} changes the base name.

\section{Native \file{.res} files}

The \file{.res} file is organized into loadcases and fields. Each loadcase starts with a loadcase header followed by any number of field blocks.

\subsection{Loadcase blocks}

\begin{syntax}
LC <id>
\end{syntax}

The loadcase ID follows execution order. The following field blocks belong to that loadcase until the next loadcase header appears.

\subsection{Dense fields}

\begin{syntax}
FIELD, NAME=<name>, TYPE=<type>, COLS=<cols>, ROWS=<rows>
<v11> <v12> ...
END FIELD
\end{syntax}

\key{TYPE} is always written and is one of \val{NODE}, \val{ELEMENT}, \val{ELEMENT_NODAL}, \val{ELEMENT_IP}, \val{ELEMENT_MP}, or \val{UNKNOWN}. \key{ROWS} gives the number of rows. \key{COLS} gives the number of values per row. For nodal fields, a row typically corresponds to a node-position row. For element fields, a row corresponds to an element-position row.

\subsection{Indexed fields}

\begin{syntax}
FIELD, NAME=<name>, TYPE=<type>, INDEX_COLS=<i>, VALUE_COLS=<v>, ROWS=<rows>
<index...> <value...>
END FIELD
\end{syntax}

Element-local fields are written with explicit location indices. For \val{ELEMENT_NODAL}, \val{ELEMENT_IP}, and \val{ELEMENT_MP} fields, the native writer uses two index columns: the element ID and the zero-based local row number within that element. These are followed by the field values.

\subsection{Common field names}

\begin{longtable}{@{}p{0.28\linewidth}p{0.62\linewidth}@{}}
\toprule
\textbf{Loadcase} & \textbf{Typical fields} \\
\midrule
\endhead
Linear Static & \val{DISPLACEMENT}, \val{STRAIN}, \val{STRESS}, \val{STRESS_TOP}, \val{STRESS_BOT}, \val{SHELL_RESULTANTS}, \val{EXTERNAL_FORCES}, \val{REACTION_FORCES}, \val{LOCAL_SECTION_FORCES}, and optional \val{SHEAR_FLOW}. \\
Linear Static Topology & Static fields plus \val{COMPLIANCE}, \val{DENS_GRAD}, \val{VOLUME}, \val{DENSITY}, optional \val{SHEAR_FLOW}, and optional \val{ORIENTATION_GRAD} and \val{ORIENTATION}. \\
Nonlinear Static & Per accepted increment: \val{DISPLACEMENT_i}, \val{STRESS_i}, \val{LAMBDA_i}. Final state: \val{DISPLACEMENT}, \val{STRAIN}, \val{STRESS}, \val{STRESS_TOP}, \val{STRESS_BOT}, \val{EXTERNAL_FORCES}, \val{INTERNAL_FORCES}, and \val{REACTION_FORCES}. \\
Eigenfrequency & \val{MODE_SHAPE_i}, \val{PARTICIPATION_i}, \val{EIGENVALUES}, \val{EIGENFREQUENCIES}, \val{FREQUENCIES}. \\
Linear Buckling & \val{BUCKLING_MODE_i}, \val{BUCKLING_FACTORS}. \\
Linear Transient & \val{DISPLACEMENT_k}, \val{VELOCITY_k}, \val{ACCELERATION_k} for written time steps. \\
Linear Harmonic & For every requested frequency: \val{DISPLACEMENT_REAL}, \val{DISPLACEMENT_IMAG}, \val{STRESS_REAL}, \val{STRESS_IMAG}, \val{STRAIN_REAL}, and \val{STRAIN_IMAG}. \\
\bottomrule
\end{longtable}

A six-component nodal motion field uses \dofvec. A six-component stress or strain field uses the engineering order \(x,y,z,yz,zx,xy\). These conventions must not be confused.

\subsection{Linear harmonic result fields}

A \val{LINEARHARMONIC} loadcase writes the complex steady-state response as separate real and imaginary nodal fields for every requested excitation frequency. If a harmonic response quantity is written as
\[
\hat q = q_\mathrm{real} + i q_\mathrm{imag},
\]
then the corresponding \val{*_REAL} and \val{*_IMAG} fields contain these two components of the complex amplitude.

\begin{longtable}{@{}p{0.29\linewidth}p{0.13\linewidth}p{0.48\linewidth}@{}}
\toprule
\textbf{Field} & \textbf{Domain} & \textbf{Components} \\
\midrule
\endhead
\val{DISPLACEMENT_REAL} & \val{NODE} & Real part of the harmonic nodal displacement amplitude in the mechanical nodal order \dofvec. \\
\val{DISPLACEMENT_IMAG} & \val{NODE} & Imaginary part of the harmonic nodal displacement amplitude in the mechanical nodal order \dofvec. \\
\val{STRESS_REAL} & \val{NODE} & Real part of the recovered harmonic nodal stress amplitude in engineering order \(x,y,z,yz,zx,xy\). \\
\val{STRESS_IMAG} & \val{NODE} & Imaginary part of the recovered harmonic nodal stress amplitude in engineering order \(x,y,z,yz,zx,xy\). \\
\val{STRAIN_REAL} & \val{NODE} & Real part of the recovered harmonic nodal strain amplitude in engineering order \(x,y,z,yz,zx,xy\). \\
\val{STRAIN_IMAG} & \val{NODE} & Imaginary part of the recovered harmonic nodal strain amplitude in engineering order \(x,y,z,yz,zx,xy\). \\
\bottomrule
\end{longtable}

For any scalar component, amplitude and phase can be reconstructed from the stored real and imaginary values as
\[
|\hat q|=\sqrt{q_\mathrm{real}^2+q_\mathrm{imag}^2},
\qquad
\varphi=\operatorname{atan2}(q_\mathrm{imag},q_\mathrm{real}) .
\]
FEMaster currently writes the real/imaginary field blocks in frequency-sweep order, but the native \file{.res} writer does not encode the supplied frequency frame value in the field header. The associated excitation frequency is therefore determined from the \cmd{FREQUENCIES} sequence and the block order.

\section{FRD files}

The \file{.frd} writer emits the model mesh once and then writes supported nodal result blocks. It is intended for visualization in CalculiX/CGX-compatible tools. Element-local fields, integration-point fields, material-point fields, participation factors, eigenvalue tables, and scalar history fields remain available in the native \file{.res} file but are not written as FRD result blocks.

FRD node and element numbers are normally the internal IDs. If the mesh contains ID 0, the FRD writer shifts those IDs by one because FRD consumers generally expect positive numbering.

Supported structural elements are written with FRD element type codes for common solid, shell, beam, and truss topologies. Higher-order wedge and hexahedral connectivity is reordered to match the FRD convention before it is written.

\begin{longtable}{@{}p{0.28\linewidth}p{0.24\linewidth}p{0.38\linewidth}@{}}
\toprule
\textbf{FEMaster field} & \textbf{FRD block} & \textbf{Components} \\
\midrule
\endhead
\val{DISPLACEMENT}, \val{MODE_SHAPE_i}, \val{BUCKLING_MODE_i} & \val{DISP} & \val{D1}, \val{D2}, \val{D3}, \val{D4}, \val{D5}, \val{D6}, \val{ALL}. \\
\val{VELOCITY_i} & \val{VELO} & \val{V1}, \val{V2}, \val{V3}, \val{ALL}. \\
\val{ACCELERATION_i} & \val{ACCE} & \val{A1}, \val{A2}, \val{A3}, \val{ALL}. \\
\val{REACTION_FORCES} & \val{FORC} & \val{F1}, \val{F2}, \val{F3}, \val{F4}, \val{F5}, \val{F6}, \val{ALL}. \\
\val{EXTERNAL_FORCES} & \val{EXTFORC} & \val{F1}, \val{F2}, \val{F3}, \val{F4}, \val{F5}, \val{F6}, \val{ALL}. \\
\val{INTERNAL_FORCES} & \val{INTFORC} & \val{F1}, \val{F2}, \val{F3}, \val{F4}, \val{F5}, \val{F6}, \val{ALL}. \\
\val{STRESS} & \val{STRESS} & \val{SXX}, \val{SYY}, \val{SZZ}, \val{SYZ}, \val{SZX}, \val{SXY}. \\
\val{STRAIN} & \val{TOSTRAIN} & \val{EXX}, \val{EYY}, \val{EZZ}, \val{EYZ}, \val{EZX}, \val{EXY}. \\
\val{MECHANICAL_STRAIN} & \val{MESTRAIN} & \val{MEXX}, \val{MEYY}, \val{MEZZ}, \val{MEYZ}, \val{MEZX}, \val{MEXY}. \\
\val{STRESS_TOP} & \val{STRPOS} & \val{SXX}, \val{SYY}, \val{SZZ}, \val{SYZ}, \val{SZX}, \val{SXY}. \\
\val{STRESS_BOT} & \val{STRNEG} & \val{SXX}, \val{SYY}, \val{SZZ}, \val{SYZ}, \val{SZX}, \val{SXY}. \\
\val{SHELL_RESULTANTS} & \val{SHR} & \val{NXX}, \val{NYY}, \val{NXY}, \val{MXX}, \val{MYY}, \val{MXY}, \val{QX}, \val{QY}. \\
\bottomrule
\end{longtable}

The harmonic names \val{DISPLACEMENT_REAL}, \val{DISPLACEMENT_IMAG}, \val{STRESS_REAL}, \val{STRESS_IMAG}, \val{STRAIN_REAL}, and \val{STRAIN_IMAG} are not currently registered as FRD field aliases. Use the native \file{.res} output for the complete harmonic real/imaginary response until dedicated FRD mappings are added.

For shell resultants, an explicitly assigned \cmd{SHELLSECTION} orientation also defines the output basis. Without an orientation, FEMaster writes physical shell stresses in global Cartesian components and rotates shell resultants into a deterministic tangent basis obtained by projecting global \(X\) into the shell plane, falling back to global \(Y\) when necessary.

For transient results, the frame value written to FRD is the physical time passed by the transient loadcase. For eigenfrequency and buckling results, the frame value is the frequency or buckling factor associated with the written mode.
